\subsection{结构弱化}

我们假设 \(\mathbf{K} = \mathbf{R}\)。设 \(r' \in \mathrm{N}_{\mathbb{K}}\) 满足 \(r' \leqslant r\)。设 X 是一个 \(\mathbf{C}^{r}\) 类流形，\(\mathbf{X}'\) 是通过结构弱化（5.13.1）得到的 \(\mathbf{C}^{r'}\) 类流形。设 \(x \in \mathbf{X}\)，并设 \(t \in \mathrm{T}_x^{(k)}(\mathbf{X})\)，其中 \(k \leqslant r'\)。取 X 的一个卡 \(c = (\mathrm{U}, \varphi, \mathrm{E})\)，其中心为 \(x\)，并用 \(\tilde{t}\) 表示 \(\mathrm{TS}^{(k)}(\mathrm{E})\) 中满足 \(\theta_c(\tilde{t}) = t\) 的元素。由于 \(c\) 也是 \(\mathbf{X}'\) 的以 \(x\) 为中心的卡，元素 \(t' = \theta_c(\tilde{t})\)（\(\mathrm{T}_x^{(k)}(\mathrm{X}')\) 中的元素）有定义；它与 \(c\) 的选取无关。映射 \(t \mapsto t'\) 是 \(\mathrm{T}_x^{(k)}(\mathbf{X})\) 到 \(\mathrm{T}_x^{(k)}(\mathbf{X}')\) 上的一个双射，借助它把这两个空间恒同。这样得到的恒同与前面各条中的运算 \(\langle f, t \rangle\)、\(\varphi_*(t)\)、\(t_1 \otimes t_2\)、\(c(t)\)、\ldots{} 相容。

